\section{Relevant operators and boundary conditions}\label{sec3}
The operator of Laplace type is not necessarily the scalar
Laplacian. In fact, in almost all models of quantum field
theory the one loop effective action is defined by an operator
of this type. This can be demonstrated by bringing relevant
operators to the canonical form (\ref{laplaceb}). 
In this section we give explicit
construction of the connection $\omega$ and the endomorphism
(matrix-valued potential)
$E$ for scalar, spinor, vector and graviton fields. We also
describe appropriate boundary conditions.

%%%%%%%%%%%%%%%
\subsection{Scalar fields}\label{s3scalar}
Consider first the example of the multi-component real scalar
field $\Phi^A$ in $n$ dimensions. The action reads
\begin{equation}
\mathcal{L}
=\int_M d^nx\, \sqrt{g} (g^{\mu\nu}\nabla_\mu\Phi^A\nabla_\nu\Phi^A
+U(\Phi ) +\xi R \Phi^A\Phi^A ) \, ,\label{scalaraction}
\end{equation}
where $\xi$ is the conformal coupling parameter, $U$ is a potential.
The covariant derivative $\nabla_\mu$ contains the background
gauge field $G_\mu^{AB}$: $\nabla_\mu\Phi^A=\partial_\mu\Phi^A+
G_\mu^{AB}\Phi^B$. $G_\mu$ is antisymmetric in internal
indices $A,\,B$.
To evaluate the one-loop effective action one should expand 
(\ref{scalaraction}) around a
background field $\bar \Phi$, $\Phi=\bar \Phi +\phi$, and keep the
terms quadratic in fluctuations:
\begin{eqnarray}
&&\mathcal{L}_2
=\int_M d^nx \sqrt{g} \phi^A\left(-(\nabla^\mu\nabla_\mu )^{AB}
+\frac 12(U(\bar\Phi )'')^{AB} +\xi R \delta^{AB} \right)
\phi^B \nonumber \\
&&\qquad -\int_{\partial M}d^{n-1}x\,\sqrt{h} \phi^A \nabla_n \phi^A \,,
\label{S2scalar}
\end{eqnarray}
$h$ is the determinant of the induced metric on the boundary.
The inner product for quantum fields $\phi$ reads
\begin{equation}
<\phi_1,\phi_2>=\int_M d^nx\, \sqrt{g} \phi_1^A \phi_2^A \ .
\label{inn-prod-scal}
\end{equation}

The operator $D$ is defined by the bulk part of the action
(\ref{S2scalar}):
\begin{equation}
D^{AB}:=-(\nabla^\mu\nabla_\mu )^{AB}
+\frac 12(U(\bar\Phi )'')^{AB} +\xi R \delta^{AB}\,.
\label{opDscalar}
\end{equation}
For a special choice of the parameter $\xi$:
\begin{equation}
\xi =\frac{n-2}{4(n-1)} \label{confxi}
\end{equation}
the operator $D$ (\ref{opDscalar}) is conformally covariant
(if also $U(\bar\Phi )''=0$).
To bring the operator (\ref{opDscalar}) to the canonical form
(\ref{laplaceb}) we introduce
\begin{equation}
\omega_\mu^{AB}=G_\mu^{AB}\, ,\qquad
E^{AB}=-\frac 12(U(\bar\Phi )'')^{AB} -\xi R \delta^{AB}\,.
\label{omEscal}
\end{equation}
For this case, $\Omega_{\mu\nu}$ is just the ordinary Yang-Mills
field strength.

The operator (\ref{opDscalar}) is symmetric with respect to the
inner product (\ref{inn-prod-scal})
if  the surface integral
\begin{equation}
\int_{\partial M} d^{n-1}x\,\sqrt{h} (\phi_1^A\nabla_n\phi_2^A -
\phi_2^A\nabla_n\phi_1^A ) \label{symm-cond-scal}
\end{equation}
vanishes for arbitrary $\phi_1$ and $\phi_2$ belonging to its' 
domain of definition. 
This may be achieved if one imposes either
Dirichlet 
\begin{equation}
\phi^A\oB =0 \label{Dirscal}
\end{equation}
or modified Neumann 
\begin{equation}
(\nabla_n\phi^A +S^{AB}\phi^B)\oB =0 \label{neuscal}
\end{equation}
boundary conditions. $S^{AB}$ is an arbitrary symmetric matrix.
Note, that the integral (\ref{symm-cond-scal}) vanishes also if
$S^{AB}$ is an arbitrary symmetric (differential) operator
on the boundary. 

For the Dirichlet conditions (\ref{Dirscal}) the boundary term
in (\ref{S2scalar}) vanishes automatically. To ensure absence
of the surface term for the modified Neumann conditions 
(\ref{neuscal}) one should add to (\ref{scalaraction}) an appropriate
surface action.
%%%%%%%%%%%
\subsection{Bosonic string}\label{s3string}
Our next example is the non-linear sigma model in two dimensions
described by the action
\begin{eqnarray}&&
\mathcal{L}^{[\sigma]}=\int_M d^2 x\left(
\sqrt{g} G_{AB}(X) g^{\mu\nu}
\partial_\mu X^A \partial_\nu X^B +
\epsilon^{\mu\nu} B_{AB}(X) \partial_\mu X^A \partial_\nu X^B \right)
\nonumber\\
&&\qquad\qquad +\int_{\partial M} A_BdX^B\,.
\label{actsigma}
\end{eqnarray}
From the point of view of two-dimensional world sheet the fields
$X^A(x)$ are scalars. In string theory they are interpreted as
coordinates on a $d$-dimensional target manifold with the metric $G_{AB}(X)$.
$\epsilon^{\mu\nu}$ is the Levi-Civita tensor density, 
$\epsilon^{12}=-\epsilon^{21}=1$.
$B_{AB}(X)$ is an antisymmetric tensor field on the target space.
$A_B(X)$ is the electromagnetic vector potential. The action 
(\ref{actsigma}) describes charged open stings. For simplicity,
we absorb the inverse string tension $\alpha'$ into a field redefinition.
We do not include tachyon and dilaton couplings in the bulk or on the
boundary.

Usually the term with the $B$-field gets an imaginary coefficient
in the Euclidean space. Since the physical space-time has Minkowski
signature it is not especially significant which way of continuation
to the Euclidean space has been chosen provided the results are
properly continued back to Minkowski space after the calculations.
As we will see below, real coefficient in front of the $B$-field leads
to a well-defined spectral problem. This situation is in close
analogy with the continuation rules for the axial vector
field in the spinor determinant \cite{Andrianov:1984qj,Andrianov:1984fg}.
The other way to deal with the field $B$ is to keep the coefficient
of the $B$-term in (\ref{actsigma}) imaginary at the expense
of introducing a more sophisticated conjugation operation 
\cite{Osborn:1991gm} containing the sign-reversion of the
$B$-field. The same refers to the electromagnetic potential $A_B$.

The field $X$ enters the action (\ref{actsigma}) at many places making
the background field expansion a quite cumbersome procedure. The most
economic way to arrange such an expansion and to calculate higher
derivatives of the action (\ref{actsigma}) is to introduce the geodesic
coordinates in the target space. A detailed explanation of the method 
as well as further references can be found in \cite{Braaten:1985is}.
Consider the target space geodesics defined
independently at each point of the two-dimensional world surface
and parametrised by the arc length $s$ in the target space. They satisfy
the usual geodesic equation
\begin{equation}
\frac{d^2}{ds^2} X^A(x,s) +{\gamma}_{BC}^A(X)
\frac{d}{ds} X^B(x,s) \frac{d}{ds} X^C(x,s)=0 \,,
\label{geod}
\end{equation}
where ${\gamma}_{BC}^A(X)$ is the Christoffel connection
corresponding to the target space metric $G_{AB}$.
Let us supplement the equation (\ref{geod}) by the initial
conditions
\begin{equation}
X^A(x,0)=\bar X^A(x)\,, \qquad \frac{d}{ds} X^A(x,0)=
\xi^A(x)\,, \label{incondgeo}
\end{equation}
where $\bar X$ is the background field. $\xi^A$ parametrises
deviations from $\bar X$ and, therefore, can be identified with
quantum fluctuations. The
$k$-th order term of the expansion of the action
(\ref{actsigma}) around the background $\bar X$ is given by
\begin{equation}
\mathcal{L}_k=\frac 1{k!} \frac{d^k}{ds^k} \mathcal{L}(X(s)) \vert_{s=0} \,.
\label{Snsig}
\end{equation}
Higher order derivatives of $X(x,s)$ with respect to $s$
can be traded for the first derivatives by means of the
geodesic equation (\ref{geod}) and then replaced by $\xi$
at $s=0$ with the help of the initial conditions (\ref{incondgeo}).

It is convenient to introduce the Riemann curvature
of the target space metric $G_{AB}$, $\mathcal{R}_{ABCD}$,
and a 3-index field strength
\begin{equation}
H_{ABC}=\partial_A B_{BC} +\partial_B B_{CA}
+\partial_C B_{BA} \,.\label{HABC}
\end{equation}
The covariant derivative $\nabla$ is
\begin{equation}
\nabla_\mu \xi^A=\partial_\mu \xi^A +\gamma_{BC}^A(\bar X)
(\partial_\mu {\bar X}^C )\xi^B +\frac 12 {\epsilon_\mu}^\nu
(\partial_\nu {\bar X}^B ){H^A}_{BC}\xi^C \,.
\label{covdsigma}
\end{equation}
The quadratic part of the action (\ref{actsigma}) reads:
\begin{eqnarray}
&&\mathcal{L}_2=\int_M d^2x \sqrt{h} \left( G_{AB}(\bar X)
 \nabla_\mu \xi^A \nabla^\mu \xi^B -\mathcal{R}_{ABCD}
(\partial_\mu \bar X^A)(\partial^\mu \bar X^D) \xi^B\xi^C \right.
\nonumber \\
&&\qquad -\frac 12 (\partial_\mu \bar X^C)(\partial_\nu \bar X^D)
\epsilon^{\mu\nu} \xi^A\xi^B 
D_AH_{BCD} \nonumber \\
&&\qquad \left. +\frac 14 (\partial_\mu \bar X^B)(\partial^\mu \bar X^D )
H_{ABC}{H^C}_{DE}\xi^A\xi^E \right) \label{S2sig} \\
&&\qquad +\int_{\partial M} d\tau \left(
D_\tau \xi^A \xi^B (B_{AB}-F_{AB}) + (\partial_\tau \bar X^B)
D_A (B_{BC}-F_{BC}) \xi^A\xi^C \right) \,,\nonumber
\end{eqnarray}
where $\tau$ is the arc length along the boundary. 
$F_{BC}=\partial_B A_C-\partial_CA_B$. The covariant derivatives
$D_\tau$ and $D_A$ contain the Christoffel connection on the target
space (but not $H_{ABC}$ as 
the full covariant derivative $\nabla$, eq. (\ref{covdsigma})).

The natural inner product in the space of fluctuations $\xi$
reads
\begin{equation}
<\xi_{(1)},\xi_{(2)}>=\int_M d^2x \sqrt{g} G_{AB}(\bar X(x))
\xi_{(1)}^A (x) \xi_{(2)}^B(x) \,.\label{inprxi}
\end{equation}
The volume part of the action (\ref{S2sig}) has now the canonical
form $\langle \xi ,D\xi \rangle$ with the operator $D$ (\ref{laplaceb})
which is obviously of Laplace type.  The connection $\omega$ is defined in 
(\ref{covdsigma}) and
the endomorphism $E$ can be easily extracted from (\ref{S2sig}).
The operator $D$ 
is formally self-adjoint with respect to (\ref{inprxi}) if
we impose the boundary conditions of Neumann type
\begin{equation}
\mathcal{B}\xi =(\nabla_n \xi^A +\mathcal{S}_B^A\xi^B)\oB =0 \label{sigbc1}
\end{equation}
with arbitrary operator $\mathcal{S}$ which should be symmetric
with respect to the restriction of (\ref{inprxi}) to the boundary:
\begin{equation}
\int_{\partial M} d\tau (\xi_{(1)}^A \mathcal{S}_A^B \xi_{(2)B}
-\xi_{(2)}^A \mathcal{S}_A^B \xi_{(1)B}) =0
\end{equation}
There is a preferable choice of the boundary operator.
Let us vary the action (\ref{S2sig}) with respect to
the fluctuation field $\xi$:
\begin{equation}
\frac 12 \delta S_2 =\int_Md^2x\, \sqrt{h} (\delta\xi ) D\xi -
\int_{\partial M} d\tau (\delta \xi ) \mathcal{B}\xi \,.
\label{varS2sig}
\end{equation}
Now we require that the boundary integral in (\ref{varS2sig}) vanishes
for arbitrary $\delta \xi$. Hence we arrive at the boundary conditions
(\ref{sigbc1}) with the operator $\mathcal{S}$ given by
\begin{eqnarray}
&&\mathcal{S}_{A}^B=\frac 12 (\Gamma \nabla_\tau +
\nabla_\tau \Gamma )_A^B +S_A^B \,,\qquad\Gamma_{A}^B={B_A}^B-{F_A}^B \,, 
\label{bopsigma} \\
&& S_A^B=\frac 14 (\partial_n \bar X^C)
\left[ {H^{DA}}_{C}(B_{DB}-F_{DB}) 
+{H^{D}}_{BC}({B_D}^A-{F_D}^A)\right]\nonumber\\
&&\qquad\qquad\qquad
+\frac 12 (\partial_\tau \bar X^C)\left[ D^A (B_{BC}-F_{BC} )
+D_B ({B^A}_C-{F^A}_C)
\right].  \nonumber 
\end{eqnarray}
Note that the operator in (\ref{bopsigma}) is not of the
ordinary Neumann (or Robin) type since it contains tangential
derivatives on the boundary (cf. sec.\ \ref{s5tang}).

The variation (\ref{varS2sig}) vanishes also if we choose Dirichlet
boundary conditions for some of the 
coordinates of the string endpoints. Namely, we can take a projector
$\Pi_+ $ and impose (\ref{sigbc1}) on $(\Pi_+)_B^A \xi^b$
and  $(1-\Pi_+)_B^A \delta\xi^B\oB =0$. 
Physically this means that the endpoints of the bosonic string
are confined to a submanifold in the target space.
Such configurations \cite{Dai:1989ua}
are called the Dirichlet branes.
%%%%%%%%%%%%%%%%%%%%%%%%%%%%%%%%%%%%%%%%
\subsection{Spinor fields}\label{s3spin}
The action for the spinor fields $\psi$ 
\begin{equation}
\mathcal{L}=\int_M d^n x \sqrt{g} \bar \psi\, \Dir\, \psi
\label{spinS}
\end{equation}
contains a first order operator $\Dir$ of Dirac type.
In Euclidean space the conjugate spinor $\bar \psi$
is just the hermitian conjugate of $\psi$: $\bar \psi=
\psi^\dag$.
By definition, an operator $\Dir$ is of Dirac type if
its square $D=\Dir^2$ is of Laplace type.
Spectral theory of general operators of Dirac type both
on manifolds without boundaries and with local boundary
conditions on manifolds with boundaries can be found
in \cite{Branson:1992a,Branson:1992b}. Here we consider
some physically motivated particular cases only. Let us introduce
the Euclidean Dirac $\gamma$-matrices which satisfy the Clifford
commutation relations:
\begin{equation}
\gamma_\mu\gamma_\nu + \gamma_\nu \gamma_\nu = 2g_{\mu\nu} \,.
\label{Clifford}
\end{equation}
The $\gamma$-matrices defined in this way are hermitian,
$\gamma_\mu^\dag =\gamma_\mu$. We also need the chirality 
matrix which will be denoted $\gamma^5$ independently of
the dimension. It satisfies
\begin{equation}
(\gamma^5)^\dag =\gamma^5 ,\qquad \gamma^5\gamma_\mu =
-\gamma_\mu \gamma^5 \,.\label{gamma5}
\end{equation}
From now on we suppose that the dimension $n$ is even. We fix the sign in
$\gamma^5$ by choosing
\begin{equation}
\gamma^5=\frac{i^{\frac{n(n-1)}2}}{n!}
\epsilon^{\mu\dots\nu}\gamma_\mu \dots \gamma_\nu \,.
\label{sign5}
\end{equation}
Let $\Dir$ be the standard Dirac operator in curved space
with gauge and axial gauge connection
\begin{equation}
\Dir =i\gamma^\mu \left(\partial_\mu +\frac 18 [\gamma_\nu ,\gamma_\rho ]
\sigma_\mu^{\nu\rho} +A_\mu +iA_\mu^5 \gamma^5 \right) \,.
\label{dirop}
\end{equation}
Here $\sigma_\mu^{\nu\rho}$ is the spin-connection (\ref{spincon}). 
$A_\mu$ and
$A_\mu^5$ are vector and axial vector fields respectively
taken in some representation of the gauge group. Both
$A_\mu$ and $A_\mu^5$ are antihermitian in the gauge indices.
The operator $\Dir$ (\ref{dirop}) is formally self-adjoint
in the bulk.
The operator $D=\Dir^2$ is of Laplace type (\ref{laplaceb}) 
with\footnote{The present author is grateful to Valery Marachevsky
for his help in deriving and checking eqs.\ (\ref{oEdir}),
(\ref{Omdir}) and (\ref{spinbc4}), see also \cite{DeBerredo-Peixoto:2001qm}.}
\begin{eqnarray}
&&\omega_\mu =\frac 18 [\gamma_\nu ,\gamma_\rho ]
\sigma_\mu^{\nu\rho} +A_\mu +\frac i2 [\gamma_\mu ,\gamma_\nu ]
A^{5\nu}\gamma^5 \,, \nonumber \\
&&E=-\frac 14 R +\frac 14 [\gamma^\mu ,\gamma^\nu ]F_{\mu\nu}
+i\gamma^5 D^\mu A_\mu^5 -(n-2)A_\mu^5A^{5\mu} \nonumber \\
&&\qquad\qquad\qquad
-\frac 14 (n-3) [\gamma^\mu ,\gamma^\nu ] [A_\mu^5 ,A_\nu^5 ]
\label{oEdir}
\end{eqnarray}
with obvious notations $F_{\mu\nu}=\partial_\mu A_\nu -
\partial_\nu A_\mu +[A_\mu ,A_\nu ]$, 
$D_\mu A_\nu^5 = \partial_\mu A_\nu^5 -
\Gamma_{\mu\nu}^\rho A_\rho^5 +[A_\mu ,A_\nu^5 ]$.
The expression for $\Omega_{\mu\nu}$ is a little bit
lengthy:
\begin{eqnarray}
&&\Omega_{\mu\nu}=F_{\mu\nu}-[A_\mu^5 ,A_\nu^5 ]-\frac 14
\gamma^\sigma \gamma^\rho R_{\sigma\rho\mu\nu} 
-i\gamma^5\gamma^\rho (\gamma_\nu D_\mu A_\rho^5
-\gamma_\mu D_\nu A_\rho^5 ) \nonumber \\
&&\qquad\quad +i\gamma^5 A_{\mu\nu}^5 
+[A_\mu^5,A_\rho^5] \gamma^\rho \gamma_\nu -
[A_\nu^5,A_\rho^5] \gamma^\rho \gamma_\mu \nonumber \\
&&\qquad\quad 
-\gamma^\rho A_\rho^5 \gamma_\mu \gamma^\sigma A_\sigma^5 \gamma_\nu
+ \gamma^\rho A_\rho^5 \gamma_\nu \gamma^\sigma A_\sigma^5 \gamma_\mu
\,,\label{Omdir}
\end{eqnarray}
where
\begin{equation}
A_{\mu\nu}^5=\partial_\mu A_\nu^5-\partial_\nu A_\mu^5
+[A_\mu ,A_\nu^5 ]-[A_\nu ,A_\mu^5 ] \label{amunu}
\end{equation}

Consider now manifolds with boundary. The specific feature of the action
(\ref{spinS}) is that it contains first order derivatives
only. Consequently, boundary conditions should be imposed
on a half of the spinor components. Let these be Dirichlet
boundary conditions
\begin{equation}
\Pi_-\psi \oB =0 \label{spinbc1}
\end{equation}
where $\Pi_-$ is a hermitian projector, $\Pi_-^2=\Pi_-$, 
$\Pi_-^\dag =\Pi_-$. Due to the hermiticity
\begin{equation}
\bar\psi\Pi_- \oB =0 \label{spinbc2}
\end{equation}
Following Luckock \cite{Luckock:1990xr}
let us consider a family of the projectors
\begin{equation}
\Pi_-=\frac 12 \left( 1+\gamma^n \exp (iq\gamma^5) \right)\,,
\label{proluck}
\end{equation}
where $q$ is a scalar which can depend on the coordinate on
the boundary.

To make the operator
$\Dir$ formally self-adjoint including the boundary we must
require that 
\begin{equation}
\int_{\partial M} d^{n-1}x\,\sqrt{h} \bar\psi_1 \gamma_n \psi_2
=0 \label{symdir}
\end{equation}
for all $\psi_{1,2}$ satisfying the boundary conditions
(\ref{spinbc1}). Since $\psi_1$ and $\psi_2$ are arbitrary,
the projector $\Pi_-$ should satisfy
\begin{equation}
(1-\Pi_-)\gamma^n (1-\Pi_-)=0 \,.\label{cpr}
\end{equation}
This condition yields $q=\pm \pi /2$. The projector (\ref{proluck})
takes the form
\begin{equation}
\Pi_-=\frac 12 \left( 1\pm i\gamma^n \gamma^5 \right)\,.
\label{proBG}
\end{equation}

To formulate the spectral problem for the {\it second}
order operator $D=\Dir^2$ we need boundary conditions for
the second half of the spinor components. The relevant functional
space should be spanned by the eigenfunctions of the Dirac
operator $\Dir$. It is clear that on this space the functions
$\Dir\psi$ should satisfy the same boundary conditions
(\ref{spinbc1}) as the $\psi$'s themselves:
\begin{equation}
\Pi_-\Dir\,\psi \oB =0 \,.\label{spinbc3}
\end{equation}
Let us adopt the choice (\ref{proBG}) for $\Pi_-$.
By commuting $\Pi_-$ with $\Dir$ in (\ref{spinbc3}) we obtain
\begin{equation}
(\nabla_n +S)\Pi_+\psi \oB =0,\qquad
S=-\frac 12 L_{aa}\Pi_+ \,,\qquad
\Pi_+=1-\Pi_- \,.\label{spinbc4}
\end{equation}
We remind that $L_{ab}$ is extrinsic curvature of the boundary.
The boundary conditions (\ref{spinbc1}), (\ref{spinbc4})
with (\ref{proBG}) are mixed (cf.\ (\ref{mixbc})).
Spectral geometry of the Dirac operator with these boundary conditions
has been thoroughly studied by Branson and Gilkey 
\cite{Branson:1992b}.

We can generalise the boundary conditions presented above by
considering non-hermitian projectors, $\Pi_-^\dag \ne \Pi_-$.
Then the boundary condition for the conjugated spinors reads:
\begin{equation}
\bar\psi \Pi_-^\dag \oB =0\,.\label{spinbc5}
\end{equation}
Instead of (\ref{cpr}) we have the condition
\begin{equation}
(1-\Pi_-^\dag )\gamma^n (1-\Pi_-)=0 \,,\label{cpr2}
\end{equation}
which yields
\begin{equation}
\Pi_-=\frac 12 \left( 1+i\gamma^n\gamma^5 e^{r(x^a)\gamma^5} \right)
\label{prochir}
\end{equation}
with an arbitrary real function (or even with an arbitrary hermitian
matrix valued function) $r(x^a)$. 

In the Minkowski signature space the spinor conjugation
includes $\gamma^0$. Therefore, the boundary conditions
on the conjugate spinor (\ref{spinbc2}) are changed. 
Roughly speaking, to continue $\Pi_-$ to the Minkowski signature
space time one has to replace $\gamma^n\gamma^5$ by $\gamma^n$
in (\ref{proBG}) and (\ref{prochir}) and to take into account
powers of $i$ which appear in the Dirac gamma matrices.

In particle physics, interest to the
boundary conditions defined by (\ref{proBG}) is due to the 
MIT bag model
of hadrons  proposed in \cite{Chodos:1974pn,Chodos:1974je,DeGrand:1975cf}. 
This model was modified later 
\cite{Chodos:1975ix,Callan:1979bm,Brown:1979ui,
Theberge:1980ye,Rho:1983bh,Goldstone:1983tu}
to include a chiral phase on the boundary in a manner similar to
(\ref{prochir})
(for a review see \cite{Hasenfratz:1978dt}).
Renormalization of quantum field theory with the boundary
conditions defined by (\ref{proBG}) was
considered by Symanzik \cite{Symanzik:1981wd}.

As we have already mentioned above, bag boundary conditions defined
by the projector (\ref{proBG}) are a particular case of mixed boundary
conditions (\ref{mixbc}). They will be considered in detail in sec.\
\ref{s5mixed} (see also sec.\ \ref{s5cases} for further references to
calculations in a ball). Chiral bag boundary conditions defined
by (\ref{prochir}) with $r\ne 0$ are considerably more complicated
because the equation (\ref{spinbc3}) contains a mixture of normal
and tangential derivatives.
We refer to 
\cite{Hrasko:1984sj,Falomir:1990tm,Falomir:1991fy,DeFrancia:1992sf,Wipf:1995dy,
DeFrancia:1996xi,DeFrancia:1998if,Beneventano:2002im,Esposito:2002vz}
for more calculations of spectral functions in this latter case. 

The boundary conditions considered in this section are {\it local},
i.e. they treat the fields at each point of the boundary independently.
One can also define {\it global} boundary conditions for the
Dirac operator (cf. sec.\ \ref{s5spec}).

Due to the fermionic nature of the spinor field the path
integral over $\psi$ and $\bar \psi$ gives determinant
of the operator $\Dir$ to a positive power,
\begin{equation}
Z=\int \mathcal{D}\bar\psi\,\mathcal{D}\psi 
\exp (-\mathcal{L}(\bar \psi ,\psi ))=\det \Dir
\label{pathspin}
\end{equation}
where the action $\mathcal{L}$ is given by (\ref{spinS}) and
$\bar\psi$, $\psi$ are complex Dirac spinors.


%%%%%%%%%%%%%%%%%%%%%%%%%%%%%%%%%%%%%%%%%%%%%%%%%%%%
\subsection{Vector fields}\label{s3vect}
Consider the Yang--Mills action for the gauge field 
$A_\mu^\alpha$ (Greek letters from the beginning of the alphabet label
generators of the gauge group):
\begin{equation}
\mathcal{L}=\frac 14 \int_M d^nx \sqrt{g} F_{\mu\nu}^\alpha
F^{\mu\nu\alpha} ,\label{YMaction}
\end{equation}
where, as usual, $F_{\mu\nu}^\alpha =
\partial_\mu A_\nu^\alpha -\partial_\nu A_\mu^\alpha
+c^\alpha_{\beta\gamma}A_\mu^\beta A_\nu^\gamma$.
$c^\alpha_{\beta\gamma}$ are totally antisymmetric
structure constants of the
gauge group. Let us introduce the background field $B_\mu^\alpha$
by the shift $A_\mu^\alpha \to B_\mu^\alpha + A_\mu^\alpha$.
From now on $A_\mu^\alpha$ plays the role of quantum
fluctuations. The quadratic part of the action reads:
\begin{eqnarray}
&&\mathcal{L}_2=\frac 12 \int_M d^nx \sqrt{g} \left[ -A^{\alpha\nu}
\nabla^\mu\nabla_\mu A_\nu^\alpha 
+ A_\mu^\alpha \nabla^\mu\nabla^\nu A_\nu^\alpha  +
A_\nu^\alpha A_\mu^\alpha R^{\mu\nu}\right. \nonumber\\
&&\qquad\qquad\qquad\qquad
\left.+2F_{\mu\nu}^\alpha (B) c^\alpha_{\beta\gamma}A^{\beta\mu}
A^{\gamma\nu} \right]
\nonumber \\
&&\qquad +\frac 12 \int_{\partial M}d^{n-1}x \sqrt{h} \left[
A_\nu^\alpha \nabla^\nu A_n^\alpha -
A^{\alpha\nu}\nabla_n A_\nu^\alpha \right] \,.
\label{S2vector}
\end{eqnarray}
The covariant derivative $\nabla$ contains both metric
and gauge parts: $\nabla_\nu A_\mu^\alpha =
\partial_\nu A_\mu^\alpha + B_\nu^\beta c^\alpha_{\beta\gamma}A_\nu^\gamma
-\Gamma_{\nu\mu}^\rho A_\rho^\alpha$. One should impose a gauge
condition on the fluctuation $A_\mu^\alpha$.
We choose
\begin{equation}
\nabla^\mu A_\mu^\alpha = 0 \,.\label{bg-gauge}
\end{equation}
In the gauge (\ref{bg-gauge}) the bulk term in (\ref{S2vector}) defines
an elliptic operator of the Laplace type with
\begin{eqnarray}
&&(\omega_\mu )^{\alpha\rho}_{\nu\beta}=
B_\mu^\gamma c^\alpha_{\gamma\beta} \delta_\nu^\rho
-\Gamma_{\mu\nu}^\rho \delta_\beta^\alpha \,,\label{omegaYM}\\
&&(E)^{\alpha\rho}_{\nu\beta} =-R_\nu^\rho \delta_\beta^\alpha
+2 {F(B)^{\gamma\rho}}_\nu c^\gamma_{\beta\alpha} \,.   \label{EYM}
\end{eqnarray}
The field strength corresponding of the connection (\ref{omegaYM})
reads
\begin{equation}
(\Omega_{\mu\nu})^{\alpha\sigma}_{\rho\beta}=
{R^\sigma}_{\rho\mu\nu}\delta^\alpha_\beta +F(B)^\delta_{\mu\nu}
c^\alpha_{\delta\beta} \delta^\sigma_\rho \,. \label{OmegaYM}
\end{equation}
Note, that all the quantities above (\ref{omegaYM}) - (\ref{OmegaYM})
are matrix functions with both gauge and vector indices.

On a manifold without boundary
the operator $D$ with (\ref{omegaYM}) and (\ref{EYM}) is
symmetric with respect to the standard inner product in the
space of the vector fields
\begin{equation}
<A^{(1)},A^{(2)}>=\int_M \sqrt{g} d^n x\ A^{(1)\alpha\mu}
A^{(2)\alpha}_\mu \,.
\label{inprYM}
\end{equation}

The ghost operator corresponding to the gauge (\ref{bg-gauge})
is just the ordinary scalar Laplacian, $D^{gh}
=-\nabla^{[gh]\mu}\nabla^{[gh]}_\mu$, with the connection
\begin{equation}
(\omega_\mu^{[gh]})^\alpha_\gamma =B_\mu^\beta c^\alpha_{\beta\gamma}\,.
\label{ghostcon}\end{equation}

The one-loop path integral reads:
\begin{equation}
Z(B)=\det (D)^{-\frac 12}_\perp \det (D^{[gh]})^{\frac 12}
\label{p-int-g1}
\end{equation}
where the first determinant is restricted to the fields
satisfying the gauge condition (\ref{bg-gauge}). Note, that
pure gauge fields $A=\nabla \xi$ are zero modes
of the total operator $D+\nabla \nabla $ of the bulk
action (\ref{S2vector}) only on shell, i.e. when the
background field $B$ satisfies the classical equation
of motion $\nabla^\mu F(B)_{\mu\nu}^\alpha =0$. Therefore,
the path integral is gauge-independent on shell, but it
depends on the gauge choice off shell. For
example, the Feynman gauge path integral
\begin{equation}
Z_F(B)=\det (D)^{-\frac 12} \det (D^{[gh]}) \,,
\label{p-int-g2}
\end{equation}
where the first determinant is calculated on the
space of {\it all} vector field, is equal to the
$Z(B)$ defined in (\ref{p-int-g1}) only on shell.
However, physical predictions of the two path integrals
are, of course, equivalent.

The path integral (\ref{p-int-g2}) can be obtained by adding
the gauge fixing term
\begin{equation}
{\mathcal{L}}_{\rm gf}=\frac 1{2\kappa} \iM (\nabla_\mu A^\mu)^2
\label{Svgf}\end{equation}
to (\ref{S2vector}) with $\kappa=1$. The case $\kappa\ne 1$
yields a non-minimal operator on the gauge field fluctuations
(cf sec.\ \ref{s4nonmin}).

One can define ``total'' heat kernel coefficients
for the path integral in a certain gauge $\chi (A)$
\begin{equation}
a^{tot}_k =a_k(D^\chi )-2a_k( D^{[gh]}_\chi ) \,,
\label{atot}
\end{equation}
where $D^\chi$ is defined by the action (\ref{S2vector})
with the gauge fixing term $\chi^2(A)$, and the ghost
operator is $D^{[gh]}_\chi \xi =\chi (\nabla \xi )$.
Even on shell only the coefficient $a_n^{tot}$ is
gauge-independent. Only this coefficient contains
information on the one-loop divergences in a gauge-invariant
regularization like the $\zeta$-function one. 

On a manifold with boundary one should impose 
boundary conditions on gauge fields and  ghosts.
The boundary conditions should be gauge invariant.
Consider a more general set-up when we have some
quantum fields $\Phi$ and a linearised
gauge transformation $\Phi\to\Phi +\delta_\xi \Phi$
with local parameter $\xi$. Boundary operator 
$\mathcal{B}$ defines gauge invariant boundary conditions
\begin{equation}
\mathcal{B}\Phi  =0
\label{ginvbc1}
\end{equation}
if there exist boundary conditions for the gauge
parameter $\xi$ 
\begin{equation}
\mathcal{B}_\xi \xi  =0 
\label{ginvbc2}
\end{equation}
such that 
\begin{equation}
\mathcal{B} \delta_\xi \Phi  =0 \,.
\label{ginvbc3}
\end{equation}
The equation (\ref{ginvbc3}) means that the
functional space defined by the boundary conditions (\ref{ginvbc1})
is invariant under the gauge transformations provided the gauge
parameter $\xi$ satisfies (\ref{ginvbc2}). Upon quantisation
(\ref{ginvbc2}) become boundary conditions for the ghost fields.
The condition (\ref{ginvbc3}) ensures validity of the 
Faddeev-Popov trick on a manifold with boundary and
guarantees gauge-independence of the on-shell path integral
\cite{Vassilevich:1998iz}.

Since the boundary term in (\ref{S2vector}) is diagonal
in the gauge index $\alpha$ we can consider the case
of an abelian gauge group and drop $\alpha$ from the
notations. The general non-abelian case does not offer
considerable complications. There are two sets of local
boundary conditions which satisfy the gauge-invariance
requirements above. The first set is called absolute
boundary conditions. It reads
\begin{eqnarray}
&&A_n\oB =0\,,\qquad \partial_n A_a \oB =
(\nabla_n \delta_{ab} -L_{ab} )A_b\oB =0 \,,
\nonumber \\
&&\nabla_n \xi \oB =0 \,.\label{absbc}
\end{eqnarray}
The second set
\begin{eqnarray}
&&(\nabla_n -L_{aa})A_n \oB =0\,, \qquad A_a\oB =0\,,
\nonumber \\
&& \xi \oB =0 \label{relbc}
\end{eqnarray}
is called relative boundary conditions. 
The projectors on the Dirichlet and Neumann subspaces and the 
endomorphism $S$ are $(\Pi^-)_{ij}=\delta_{in}\delta_{jn}$,
$(\Pi^+)_{ij}=\delta_{ij}-\delta_{in}\delta_{jn}$, 
$S_{ij}=-L_{ab}\delta_{ia}\delta_{jb}$ and
$(\Pi^-)_{ij}=\delta_{ij}-\delta_{in}\delta_{jn}$,
$(\Pi^+)_{ij}=\delta_{in}\delta_{jn}$,
$S_{ij}=-L_{aa}\delta_{in}\delta_{jn}$ for the absolute and
relative boundary conditions respectively.

It is straightforward
to check gauge invariance of the absolute boundary conditions
(\ref{absbc}). A bit more job is required to show that the
condition (\ref{relbc}) for the normal component is also gauge
invariant. The key observation is that near the boundary
the scalar Laplacian can be represented as
$D^{[gh]}=-(\nabla_n -L_{aa})\nabla_n +\mathcal{E}$,
where $\mathcal{E}$ does not contain normal derivatives
and, therefore, $\mathcal{E}\xi \oB =0$ under the boundary
conditions (\ref{relbc}) for the $\xi$. Consequently,
\begin{equation}
(\nabla_n -L_{aa}) \delta_\xi A_n \oB =
(\nabla_n -L_{aa})\nabla_n \xi \oB =
(-D^{[gh]}\xi +\mathcal{E}\xi )\oB=0 
\end{equation}
where $D^{[gh]}\xi \oB=0$ on the eigenfunctions of the
operator $D^{[gh]}$ which is enough for our purposes.

The boundary term in the action (\ref{S2vector}) can be
rewritten as\footnote{The equation to follow contains two types
of covariant derivatives. Definitions can be found in sec.\ \ref{s2geo}.}
\begin{equation}
\frac 12 \int_{\partial M} d^{n-1}x \sqrt h 
[A_a(A_{n:a} +L_{ab}A_b -A_{a;n} )] \,.
\label{btvect}
\end{equation}
This boundary action vanishes for both types of
the boundary conditions (\ref{absbc}) and (\ref{relbc}).
Hence the operator $D$ is symmetric for these
boundary conditions. Another remarkable fact is that
the fields $A^\perp$ satisfying the gauge condition
(\ref{bg-gauge}) are orthogonal to the gauge transformations.
Indeed,
\begin{equation}
<A^\perp ,\nabla \xi > =-\int_{\partial M} d^{n-1}x \sqrt h
A^\perp_n \xi =0
\end{equation}
for both (\ref{absbc}) and (\ref{relbc}). 

We should also check whether the gauge condition (\ref{bg-gauge})
is indeed compatible with the boundary conditions (\ref{relbc})
and (\ref{absbc}). For an arbitrary $A$ there should exist
unique $\xi$ such that
\begin{equation}
\nabla^\mu (A_\mu + \nabla_\mu \xi )=0 .\label{compatib}
\end{equation}
We rewrite (\ref{compatib}) as
\begin{equation}
\nabla^\mu A_\mu =-\Delta \xi \,. \label{compatib2}
\end{equation}
Let us start with relative boundary conditions (\ref{relbc}).
In this case the left hand side of (\ref{compatib2}) satisfies
Dirichlet boundary conditions. The scalar Laplacian
$\Delta$ for these boundary conditions is invertible. 
Therefore, a solution for 
(\ref{compatib2}) always exists, and it is unique. The case of
absolute boundary conditions is a bit more involved.
One should take care of a zero mode in the ghost sector.
We leave this case as an exercise. A more extensive 
discussion of compatibility of gauge and boundary conditions
can be found in \cite{Vassilevich:1998iz}.

%%%%%%%%%%%%%%%%%%%%%%%%%%%%%%%%%%%%%%%%%%%%%%%%%%%%%%%%%%%
\subsection{Graviton}
We start this section with the Einstein-Hilbert action
on a four dimensional Euclidean manifold without boundary:
\begin{equation}
\mathcal{L}=\frac 1{16\pi G_N} \int d^4 x\sqrt{g} (R-2\Lambda )\,,
\label{gract}
\end{equation}
where $R$ is the scalar curvature,
$G_N$ is the Newton constant, $\Lambda$ is the
cosmological constant.
 As usual, let us shift the metric 
$g_{\mu\nu}\to g_{\mu\nu} +h_{\mu\nu}$. From now on
$g_{\mu\nu}$ will denote the background metric, $h_{\mu\nu}$
will be the quantum fluctuations. We can decompose the 
$h_{\mu\nu}$ further in trace, longitudinal and transverse-traceless
part:
\begin{equation}
h_{\mu\nu} =\frac 14 hg_{\mu\nu} +(L\xi )_{\mu\nu} +
h^\perp_{\mu\nu} \,, \label{DYdec}
\end{equation}
where $g^{\mu\nu}h^\perp_{\mu\nu}=\nabla^\nu h^\perp_{\mu\nu}=0$
and
\begin{equation}
(L\xi )_{\mu\nu}=\nabla_\mu \xi_\nu +\nabla_\nu \xi_\nu
-\frac 12 g_{\mu\nu} \nabla^\rho \xi_\rho \,.
\label{opL}
\end{equation}
The decomposition (\ref{DYdec}) is orthogonal with respect to
the inner product:
\begin{eqnarray}
&& <h, h'>=\int d^4x \sqrt {g}
G^{\mu \nu \rho \sigma} h_{\mu \nu}
h'_{\rho \sigma} \,,\label{innprgrav} \\
&&G^{\mu \nu \rho \sigma}=\frac 12 (g^{\mu \rho} g^{\nu \sigma}+
g^{\mu \sigma} g^{\nu \rho} +Cg^{\mu \nu}g^{\rho \sigma})
\,. \nonumber 
\end{eqnarray}
Here $C$ is a constant. For positivity of (\ref{innprgrav}) $C$
should be greater than $-\frac 12$.
Under the action of infinitesimal diffeomorphism generated by
a vector $\epsilon_\mu$ the components of (\ref{DYdec}) transform
as
\begin{equation}
\xi_\mu \to \xi_\mu +\epsilon_\mu \,,\quad
h \to h+2\nabla^\mu \epsilon_\mu\,,\quad
h_{\mu \nu}^\perp \to h_{\mu \nu}^\perp \label{diftr}
\end{equation}
One can fix the gauge freedom (\ref{diftr}) by the condition
\begin{equation}
\xi_\mu =0\,. \label{gfgr}
\end{equation}
If the background admits conformal Killing vectors (these are
the vectors which are annihilated by the operator $L$
(\ref{opL})) the condition (\ref{gfgr}) is not enough and
one should impose one more gauge condition on the trace
part (see e.g. \cite{Vassilevich:1993rk}). We suppose
that conformal Killing vectors are absent.

The Jacobian factor induced by the change of variables
(\ref{DYdec}) $h_{\mu\nu}\to (h,\xi_\mu ,h_{\mu\nu}^\perp )$
is
\begin{equation}
J=\det_V (L^\dag L)^{\frac 12} \,,\label{Jgrav}
\end{equation}
where the determinant is calculated on the space of the vector
fields (excluding the conformal Killing vectors which we do not
take into account). It is convenient to shift the scalar
part of the metric fluctuations by $\nabla^\mu \xi_\mu$ so that
the decomposition (\ref{DYdec}) becomes
\begin{equation}
h_{\mu\nu} =\frac 14 (\sigma +2\nabla^\rho \xi_\rho )
g_{\mu\nu} +(L\xi )_{\mu\nu} +
h^\perp_{\mu\nu} \,. \label{DYnew}
\end{equation}
Since the change of the variables $h\to \sigma$ does not introduce
any Jacobian factor we conclude that the path integral measure is
\begin{equation}
\mathcal{D}h_{\mu\nu} =\det_V (L^\dag L)^{\frac 12}
\mathcal{D}\sigma \mathcal{D}\xi \mathcal{D}h^\perp \,.
\label{measgr}
\end{equation}

To simplify the discussion we suppose that the background metric
$g_{\mu\nu}$ satisfies the Einstein equations
\begin{equation}
R_{\mu\nu}(g)=\Lambda g_{\mu\nu} \,. \label{Einsteq}
\end{equation}
The quadratic part of the action reads:
\begin{eqnarray}
&&\mathcal{L}_2=\frac 1{16\pi G} \int d^4x \sqrt {g} \left(\frac 14
h^{\perp \mu \nu}(-\Delta g_{\mu \rho } g_{\nu \sigma}
+2R_{\mu \rho \nu \sigma})h^{\perp \rho \sigma}\right. \nonumber \\
&&\qquad\qquad\qquad \left.
-\frac 3{32} \sigma \left(-\Delta -\frac R3 \right) \sigma \right) \,.
\label{S2grav}
\end{eqnarray}
Due to the gauge invariance (\ref{S2grav}) does not contain $\xi$.
Functional integration over $\xi$ produces an infinite constant
equivalent to the volume of the diffeomorphism group which
will be neglected.

The kinetic term for $\sigma$ has a wrong (negative) sign. This
represents the well known conformal factor problem of quantum
general relativity. Different explanations of this phenomenon
\cite{Gibbons:1978ac,Schleich:1987fm,Mazur:1990by,Novozhilov:1991tf}
suggest (roughly) the same remedy: the conformal mode must be rotated
to the imaginary axis, $\sigma\to i\sigma$.

The path integral can be written in terms of functional determinants
corresponding to vector fields (indicated by the subscript $V$),
scalars (indicated by $S$), and transverse traceless tensors
(indicated by $T\perp$): 
\begin{equation}
Z=\det_V (L^\dag L)^{\frac 12}\det_{S} \left(-\Delta -\frac R3 
\right)^{-\frac 12} \det_{T\perp}(-\Delta g_{\mu \rho } g_{\nu \sigma}
+2R_{\mu \rho \nu \sigma})^{-\frac 12}
\,.\label{Zgrav}
\end{equation}
This expression is, in principle,
suitable for calculations on some homogeneous spaces since the harmonic
expansion on such spaces usually respects separation of tensors
to transverse and longitudinal parts. However, we must warn the reader
again that we have neglected the presence 
of Killing vectors and of conformal Killing vectors. 
The ways to treat these vectors created a long discussion in the
literature 
\cite{Gibbons:1978ji,Christensen:1980iy,Fradkin:1984mq,Allen:1986ta,
Griffin:1989sq,Taylor:1990ua,Vassilevich:1993rk,Mottola:1995sj,
Volkov:2000ih}.

The representation (\ref{Zgrav}) is not convenient on generic
manifold since the tensor operator is restricted to the
transverse modes. To remove this restriction one can either
multiply (\ref{Zgrav}) by a compensating vector determinant, or
add a suitable gauge fixing term to the classical action
(\ref{gract}) and repeat the quantisation procedure right from
the beginning. The result can be found elsewhere in the
literature (see, e.g., 
\cite{Gibbons:1978ji,Christensen:1980iy,Fradkin:1984mq,Taylor:1990ua}).

Boundary conditions for one-loop Euclidean quantum gravity
must be diffeomorphism invariant and must lead to ``good''
(i.e. symmetric and elliptic) operators describing metric
and ghost fluctuations. The problem of finding a suitable set
of boundary conditions in gravity appeared to be much harder
that for the lower spin fields. There exist several
proposals on the market
\cite{Barvinsky:1987dg,Luckock:1990xr,Esposito:1995aj,Vassilevich:1995zk,
Marachevsky:1996dr,
Avramidi:1996ae,Avramidi:1997hy,Avramidi:1998sh,Moss:1997ip,Esposito:1998hp}.
However, neither of these proposals is fully satisfactory.
It is clear now that proper boundary conditions must depend on
tangential derivatives. Moreover, they probably must be non-local.

The material of this section can be generalised to the case
when an independent torsion field is present.
Explicit expressions for relevant operators acting on fields 
with different spin in the Riemann--Cartan space can be found in
\cite{Obukhov:1983mm,Cognola:1988qw,Yajima:1988eh}.


